% =====================================================================
% §4 Two-window drift quantification (the core contribution)
% =====================================================================

\section{Two-window deployment-drift quantification}
\label{sec:drift}

We run the deployed pipeline against live ISO-NE data on two 7-day
windows one year apart: $\mathrm{W1} = 2025\text{-}05\text{-}01$ to
$05\text{-}07$ and $\mathrm{W2} = 2026\text{-}04\text{-}28$ to
$05\text{-}04$. Both windows are late spring (no AC, no winter heating)
to control for season. Three trained-model variants (baseline, Chronos
zero-shot, ensemble) are compared against three naive baselines
(persistence-1d, persistence-7d, climatological mean over 4 weeks; see
Appendix~\ref{app:baselines}).

\paragraph{Headline result.}
Table~\ref{tab:drift_overall} reports overall MAPE across all six
model variants on both windows, with bootstrap 95~\% CIs over
(forecast-day, zone) pairs (procedure in Appendix~\ref{app:bootstrap}).
The trained baseline degrades from \PHWONEBASELINE{} on W1 to
\PHWTWOBASELINE{} on W2 (W2 - W1 = \PHDIFFBASELINE{}~pp), while the
Chronos-Bolt-mini zero-shot model and the ensemble drift much less:
\PHWONECHRONOS{} $\to$ \PHWTWOCHRONOS{} and
\PHWONEENSEMBLE{} $\to$ \PHWTWOENSEMBLE{}.

\begin{table}[h]
\centering\small
\caption{Overall MAPE [95~\% bootstrap CI] on the two deployment windows.
The trained baseline (CNN-Transformer with HRRR weather + per-zone
demand) loses ground year-over-year; the weather-agnostic Chronos
zero-shot variant degrades much less. Naive baselines provide a
non-trivial floor in 2026.}
\label{tab:drift_overall}
\begin{tabular}{@{}lll@{}}
\toprule
Model & $\mathrm{W1}$: 2025-05-01..07 & $\mathrm{W2}$: 2026-04-28..05-04 \\
\midrule
\multicolumn{3}{l}{\textit{Trained models}} \\
Baseline (CNN-Transformer + HRRR) & \PHWONEBASELINE{}                 & \PHWTWOBASELINE{} \\
Chronos-Bolt-mini (zero-shot)     & \PHWONECHRONOS{}                  & \PHWTWOCHRONOS{} \\
Ensemble (per-zone $\alpha$)      & \PHWONEENSEMBLE{}                 & \PHWTWOENSEMBLE{} \\
\midrule
\multicolumn{3}{l}{\textit{Naive baselines}} \\
Persistence-1d                    & \PHWONEPERSI{}                    & \PHWTWOPERSI{} \\
Persistence-7d                    & \PHWONEPERS{}                    & \PHWTWOPERS{} \\
Climatological mean (4-week)      & \PHWONECLIM{}                     & \PHWTWOCLIM{} \\
\bottomrule
\end{tabular}
\end{table}

Two observations: (i) the trained baseline's degradation between W1 and
W2 is a clean year-over-year jump that exceeds the bootstrap CI band,
and (ii) the foundation-model ensemble inherits Chronos's degradation
pattern while remaining substantially better than the lag-1d
persistence floor.

\paragraph{Daily MAPE timeline.}
Figure~\ref{fig:per_day} shows daily MAPE across both windows. The
day-to-day spread within a window is small relative to the W1 $\to$ W2
jump, reinforcing that the year-over-year shift is not a noise artifact
of a single forecast issuance time.

\begin{figure}[h]
\centering
\includegraphics[width=0.95\linewidth]{per_day_mape_timeline.png}
\caption{Daily MAPE for each model variant across both 7-day windows.
The vertical gap between W1 (left) and W2 (right) for the baseline
exceeds within-window day-to-day variability, supporting the
year-over-year drift claim with multiple repeats per window.}
\label{fig:per_day}
\end{figure}

\paragraph{Per-zone breakdown.}
Figure~\ref{fig:per_zone_drift} dissects baseline / Chronos / ensemble
MAPE by zone in both windows. The baseline's degradation is sharply
zone-localised: Maine and metro Boston (NEMA-Boston) stay near their
2022 numbers; the southern Massachusetts zones (SEMA, WCMA) and Rhode
Island swing from \PHRISEMAEONE{} to \PHRISEMAEMETWO{} between W1 and
W2. The full 8-zone $\times$ 6-model breakdown is in
Appendix~\ref{app:tables} (Table~\ref{tab:full_w1}-\ref{tab:full_w2}).

\begin{figure}[h]
\centering
\includegraphics[width=0.95\linewidth]{multi_window_per_zone_mape.png}
\caption{Per-zone MAPE for the three trained-model variants on both
windows. Solid bars = W1 (2025); hatched bars = W2 (2026). The dense
southern Massachusetts and Rhode Island zones drift dramatically; Maine
and NEMA-Boston are near-stationary.}
\label{fig:per_zone_drift}
\end{figure}
